%% Abstract variants, one per target venue.
%%
%% main.tex carries the PRIMARY version below.  To retarget the manuscript,
%% replace the body of its abstract environment with one of the others and
%% re-run scripts/../abswc against the stated limit.
%%
%% This file is reference material and is not \input by main.tex.

%% ------------------------------------------------------------------------
%% PRIMARY -- Nature Communications
%% limit 150 words, this version 149
%% ------------------------------------------------------------------------
\noindent
Software increasingly acts through chains: a person instructs an agent, which
instructs sub-agents, so the process that finally acts is several hand-offs from
the party that wanted it. Each hand-off restates the specification imperfectly
and traces responsibility back only in part. Making delegation depth a strategic
choice in an evolutionary model of a competitive artificial intelligence race,
we find the two losses dangerous only together. Alone, each lifts unsafe
behaviour by less than $0.02$; jointly they lift it to $0.32$, because
responsibility that fades with depth disables the market's own answer to
degraded instructions, which is to delegate less. In the resulting equilibrium
every principal issues the safest available instruction, the population executes
unsafe ones anyway, and depth is an unbounded shelter from liability. A
statutory floor under attributed responsibility closes it, tracing the same
safety and welfare frontier as a cap on chain length without anyone counting
hand-offs.

%% ------------------------------------------------------------------------
%% FALLBACK -- JAIR / JAAMAS
%% limit 250 words, this version 247
%% ------------------------------------------------------------------------
\noindent
Software increasingly acts through chains of delegation: a principal instructs
an agent, the agent instructs sub-agents, and the process that finally acts is
several hand-offs from the party that wanted the action taken. Governance is
written at the two ends of such a chain and treats the hand-offs as plumbing.
They are not. Each hand-off restates a specification imperfectly, so a safety
instruction survives $d$ hand-offs with probability $(1-\varepsilon)^{d}$, and
each hand-off traces responsibility back only in part, so a principal at depth
$d$ is charged the fraction $\phi^{d}$ of the harm it causes. We make delegation
depth part of the strategy in a four-strategy model of a competitive artificial
intelligence race. Depth is an unbounded liability shelter: realised harm
saturates while attributed harm decays, so for every $\phi<1$ there is a depth
beyond which each further hand-off lowers the bill and raises the damage. The
two losses are dangerous only together. Each alone lifts the long-run unsafe
frequency by less than $0.02$ and jointly they lift it to $0.32$, because
per-layer attribution disables the market's own answer to degraded instructions,
which is to shorten its chains. In equilibrium every principal issues the safest
available instruction and the population executes unsafe designs anyway.
Instruments acting on the incentive to delegate behave monotonically while
instruments acting on transmission fidelity do not. A statutory floor under
attribution traces the same safety and welfare frontier as a hard cap on chain
length, and asks no regulator to count hand-offs.

%% ------------------------------------------------------------------------
%% ORIGINAL -- technical-report version, no venue
%% limit none words, this version 445
%% ------------------------------------------------------------------------
\noindent
We study what happens to a safety instruction when it is passed down a chain of
delegated agents before anything is done with it. A principal enters a
competitive race but does not act in it: it issues a specification to an agent,
which issues one to a sub-agent, and so on for $d$ hand-offs. Two things happen
along the chain, and both are geometric in $d$. The specification is
transmitted imperfectly, so it survives $d$ hand-offs intact with probability
$(1-\varepsilon)^{d}$. Responsibility is traced back only in part, so a
principal at depth $d$ is charged the fraction $\phi^{d}$ of the harm it causes.
We embed both in the reduced four-strategy AI race of Fern\'andez Domingos and
Han, make delegation depth part of the strategy, and analyse the resulting
dynamics. Five results follow. First, depth is an unbounded liability shelter.
Realised harm saturates while attributed harm decays, so for every $\phi<1$
there is a depth beyond which each further hand-off strictly lowers the bill
even though it strictly raises the damage. No level of liability stops chains
from lengthening past it. Second, and against the intuition that the two
mechanisms simply add up, neither is dangerous on its own. Start from a race
that liability has already made safe at depth zero. Imperfect transmission alone
lifts the long-run unsafe frequency to $0.018$ and imperfect attribution alone
to $0.012$; together they lift it to $0.322$, so $91\%$ of the joint effect is
an interaction. A counterfactual that freezes the population composition
isolates its single cause. Under full pass-through the market answers
transmission loss by \emph{shortening} its chains, from six hand-offs to two,
and per-layer attribution is precisely what disables that answer. Third, the
resulting equilibrium is one in which every principal issues the safest
available instruction and the population executes the unsafe designs anyway:
declared and executed behaviour differ by $0.74$ in total variation. Fourth,
instruments that act on the incentive to delegate behave monotonically while
instruments that act on transmission fidelity do not. A partial fidelity
improvement can leave the ecosystem worse off than no improvement at all,
because over-specification was an equilibrium response to erosion. A ceiling of
one hand-off removes the unsafe behaviour entirely and raises social payoff from
$-0.8$ to $60.5$. Fifth, the shelter turns on attribution vanishing in depth
and not on its geometric form: splitting the blame equally across the chain
reproduces it exactly, while any rule bounded away from zero closes it. A
statutory floor under attribution traces the same safety and welfare frontier
as a hard cap on chain length, which matters because a floor asks no regulator
to count how many hands an instruction passed through.

